\documentclass[11pt,a4paper]{article}

\usepackage[utf8]{inputenc}
\usepackage[T1]{fontenc}
\usepackage{lmodern}
\usepackage[margin=1in]{geometry}
\usepackage{amsmath,amssymb,amsthm}
\usepackage{hyperref}
\usepackage{xcolor}
\usepackage{listings}
\usepackage{booktabs}
\usepackage{enumitem}
\usepackage{float}
\usepackage{tikz}
\usetikzlibrary{arrows.meta,positioning,shapes.geometric,calc}
\usepackage{fancyhdr}
\usepackage{titlesec}
\usepackage{tabularx}

\definecolor{modblue}{HTML}{1E3A5F}
\definecolor{modcyan}{HTML}{0EA5E9}
\definecolor{codebg}{HTML}{F8F9FA}
\definecolor{codeframe}{HTML}{DEE2E6}

\hypersetup{
    colorlinks=true,
    linkcolor=modblue,
    citecolor=modblue,
    urlcolor=modcyan,
    pdftitle={Scaling Off-Chain Open-Source Management with Merkle Trees},
    pdfauthor={MOD Protocol},
}

\lstdefinestyle{solidity}{
    backgroundcolor=\color{codebg},
    frame=single,
    rulecolor=\color{codeframe},
    basicstyle=\ttfamily\small,
    keywordstyle=\color{modblue}\bfseries,
    commentstyle=\color{gray},
    stringstyle=\color{teal},
    breaklines=true,
    numbers=left,
    numberstyle=\tiny\color{gray},
    numbersep=8pt,
    tabsize=4,
    showstringspaces=false,
    morekeywords={function,external,internal,public,private,view,pure,returns,
                  require,mapping,address,uint256,bytes32,bool,memory,calldata,
                  struct,event,emit,modifier,contract,interface,pragma,solidity},
}

\lstdefinestyle{python}{
    backgroundcolor=\color{codebg},
    frame=single,
    rulecolor=\color{codeframe},
    basicstyle=\ttfamily\small,
    keywordstyle=\color{modblue}\bfseries,
    commentstyle=\color{gray},
    stringstyle=\color{teal},
    breaklines=true,
    numbers=left,
    numberstyle=\tiny\color{gray},
    numbersep=8pt,
    tabsize=4,
    showstringspaces=false,
    language=Python,
}

\pagestyle{fancy}
\fancyhf{}
\fancyhead[L]{\small\textcolor{gray}{MOD Protocol -- Merkle-Tree Module Registry}}
\fancyhead[R]{\small\textcolor{gray}{v0.2 -- August 2026}}
\fancyfoot[C]{\thepage}
\renewcommand{\headrulewidth}{0.4pt}

\titleformat{\section}{\Large\bfseries\color{modblue}}{\thesection}{1em}{}
\titleformat{\subsection}{\large\bfseries\color{modblue!80}}{\thesubsection}{1em}{}

\newtheorem{definition}{Definition}[section]
\newtheorem{property}{Property}[section]
\newtheorem{claim}{Claim}[section]

\begin{document}

\begin{titlepage}
\centering
\vspace*{3cm}

{\Huge\bfseries\textcolor{modblue}{Scaling Off-Chain\\Open-Source Management}\par}
\vspace{0.6cm}
{\Large\textcolor{gray}{Storing the tree, not the leaves:\\a Merkle-root registry for the MOD orbit ecosystem,\\and the mod protocol it registers}\par}

\vspace{2cm}
{\large Version 0.2\par}
{\large August 2026\par}

\vspace{3cm}

\begin{abstract}
\noindent
The MOD protocol today registers each orbit module as an independent on-chain row keyed by an IPFS CID. With 300+ modules already deployed and projected growth into the tens of thousands as autonomous agents publish their own, the linear per-module storage and gas cost becomes the binding constraint on ecosystem scale. We propose replacing the per-module hash table with a single 32-byte Merkle root anchored on-chain. The full tree---containing module metadata, code CIDs, owner addresses, prices, and version history---lives entirely off-chain in IPFS-backed mirrors. Authenticity is preserved by Merkle inclusion proofs; updates are batched into a single root-rotation transaction; reads are free. We show that this design reduces per-module on-chain storage from $\approx 5$ slots ($\approx 100{,}000$ gas amortised) to $O(1/N)$, reduces a publish operation from $\approx 80{,}000$ gas to a marginal cost approaching the cost of a single SSTORE shared across $N$ updates, and makes the registry compatible with permissionless off-chain indexers without trust assumptions. The trade-off is a one-block latency between local commit and global root rotation, which is acceptable for a developer-tools registry. Part~II then specifies the mod protocol itself---what the registry registers: the \emph{module} as a directory whose anchor class \emph{is} its API; the equivalence of the Python, CLI, and HTTP surfaces and the null call that makes every module self-describing; a single addressing rule under which a name resolves to an app, an API, and a DNS record by construction; a shared Nix environment and supervisor that run every module without per-module containers; a federated gateway with scale-to-zero execution; one shared identity module issuing time-bounded signature tokens, with an owner/peer authorization model that keeps access control off-chain while the public tree carries only public facts; and the introspection, \texttt{skill.md}, and MCP surfaces that make a module callable by an agent with no per-module integration. Part~II closes with the nine invariants that constitute the protocol.
\end{abstract}

\vfill
\end{titlepage}

\tableofcontents
\newpage

\part{The Registry Protocol}

\section{Motivation}

\subsection{Current design}

The current \texttt{Registry.sol} stores one struct per module:

\begin{lstlisting}[style=solidity]
struct Module {
    address owner;
    string  name;
    string  data;   // IPFS CID
    uint256 modId;
}
mapping(uint256 => Module) public modules;
mapping(address => uint256[]) public userMods;
\end{lstlisting}

Each \texttt{registerMod} call writes:
\begin{itemize}
    \item the struct (typically 4--5 SSTOREs after string encoding overhead),
    \item an append to \texttt{userMods},
    \item a monotonic counter bump.
\end{itemize}

At 80{,}000--110{,}000 gas per registration on Base, and re-registration on every CID update, the cost is dominated by storage growth in the registry contract itself.

\subsection{The scaling wall}

Three trends make the per-row design untenable:

\begin{enumerate}
    \item \textbf{Agentic publishing.} The \texttt{agent} module already composes new modules autonomously, and a module an agent writes is callable and registrable by exactly the same rules as a hand-written one (Part~II); a 24-hour run can produce tens of new mods.
    \item \textbf{Versioning.} Most production modules update weekly. The registry currently has no compression: each version is a fresh write.
    \item \textbf{Permissionless mirrors.} Indexers that want to expose the registry over HTTP must replay every \texttt{ModRegistered} event since genesis, which is a quadratic cost in time and state.
\end{enumerate}

\subsection{Goal}

Compress the on-chain footprint to a single root, while preserving:

\begin{itemize}
    \item authenticity (no off-chain actor can lie about what's registered),
    \item revocability (owners can delete or transfer modules),
    \item ownership (only the rightful owner can update an entry),
    \item discoverability (clients can enumerate without trust).
\end{itemize}

\section{Construction}

\subsection{The off-chain tree}

\begin{definition}[Module record]
A module record $r$ is a tuple
\[
r = (\text{name},\ \text{owner},\ \text{cid},\ \text{price},\ \text{version},\ \text{updatedAt})
\]
serialised with canonical JSON (RFC 8785) and hashed with \texttt{keccak256}:
\[
\ell(r) = \mathrm{keccak256}(\mathrm{cjson}(r)).
\]
\end{definition}

\begin{definition}[Module tree]
Given $N$ module records $\{r_1, \ldots, r_N\}$, the module tree $\mathcal{T}$ is a binary Merkle tree over the multiset $\{\ell(r_1), \ldots, \ell(r_N)\}$ sorted by $\ell$ (sorted-pair construction; identical to OpenZeppelin's \texttt{MerkleProof} library). The root is
\[
R = \mathrm{root}(\mathcal{T}).
\]
\end{definition}

\subsection{On-chain anchor}

\begin{lstlisting}[style=solidity]
contract TreeRegistry {
    bytes32 public root;          // current Merkle root
    uint256 public epoch;         // rotation counter
    address public publisher;     // authorised root publisher (later DAO)
    mapping(bytes32 => bool) public revoked;  // tombstones

    event Rotated(uint256 indexed epoch, bytes32 indexed root, bytes32 cid);
    event Revoked(bytes32 indexed leaf);

    function rotate(bytes32 newRoot, bytes32 manifestCid) external {
        require(msg.sender == publisher);
        root = newRoot;
        epoch++;
        emit Rotated(epoch, newRoot, manifestCid);
    }

    function verify(bytes32 leaf, bytes32[] calldata proof)
        external view returns (bool)
    {
        if (revoked[leaf]) return false;
        return MerkleProof.verifyCalldata(proof, root, leaf);
    }

    function revoke(bytes32 leaf, bytes32[] calldata proof) external {
        // Only the owner encoded in the leaf can revoke.
        // Owner check performed off-chain by reconstructing the record
        // and verifying that msg.sender matches r.owner.
        // ...
        revoked[leaf] = true;
        emit Revoked(leaf);
    }
}
\end{lstlisting}

The on-chain state is exactly:
\begin{itemize}
    \item one \texttt{bytes32} root,
    \item one \texttt{uint256} epoch,
    \item one \texttt{address} publisher,
    \item a sparse \texttt{revoked} map (only paid for by the revoker, only for actively revoked entries).
\end{itemize}

\subsection{The publisher and the manifest}

The publisher is the actor authorised to call \texttt{rotate}. In Phase 1 it is the protocol multisig; in Phase 2 it is a threshold of validators sourced from \texttt{StakeTime}.

On every rotation the publisher uploads the full tree to IPFS as a \emph{manifest}:

\begin{lstlisting}[basicstyle=\ttfamily\small,frame=single,backgroundcolor=\color{codebg}]
{
  "epoch": 412,
  "root": "0x9f4c...e2",
  "prev": "0x2cb1...0a",
  "records": [
    { "name": "agent",     "owner": "0x...", "cid": "Qm...", ... },
    { "name": "bloctime",  "owner": "0x...", "cid": "Qm...", ... },
    ...
  ],
  "tree": { "depth": 18, "nodes": "..." }
}
\end{lstlisting}

The manifest CID is emitted in the \texttt{Rotated} event so that any indexer with an IPFS gateway can pull the new tree and verify its root against the on-chain value.

\subsection{Update flow}

\begin{enumerate}
    \item Developer signs a record update with their wallet.
    \item Signed update is broadcast to a public mempool (a Caddy-fronted FastAPI endpoint at \texttt{/api/whitepaper/tree/submit}).
    \item Publisher (or any node) aggregates updates between epochs.
    \item Publisher recomputes the tree, uploads the manifest, calls \texttt{rotate(newRoot, manifestCid)}.
    \item Clients verify their relevant leaves against the new root via Merkle proofs.
\end{enumerate}

A single \texttt{rotate} transaction subsumes an arbitrary number of registrations and updates. The cost per individual update is $O(1/N)$ where $N$ is the rotation batch size.

\section{Chain Anchor: StakeTime as the Publisher}
\label{sec:chain}

A Merkle-tree registry shifts cost off chain, but raises a new question: \emph{who is trusted to rotate the root?} We answer this with \textbf{StakeTime}, the time-weighted staking and validator-consensus primitive deployed on Base.

\subsection{StakeTime primitives}

StakeTime composes five independent on-chain primitives:

\begin{itemize}
    \item \textbf{Staking} --- delegated, lock-weighted positions that mint synthetic STT proportional to $\text{amount} \cdot M(\text{lockBlocks})$.
    \item \textbf{Consensus} --- pluggable scoring (Yuma decay, linear, stake-weighted, privacy-preserving).
    \item \textbf{Inflation} --- pluggable emission curve (halving, flat, linear decay, sigmoid).
    \item \textbf{Subnet} --- emission-token-bearing partitions of the validator set.
    \item \textbf{Registry} --- competitive 420-slot directory with weakest-replacement when full.
\end{itemize}

\subsection{Publisher election}

The publisher role is taken by any validator selected by the active consensus adapter. Their STT-weighted score $\sigma_v$ determines:

\begin{enumerate}
    \item \textbf{Right to rotate}: only validators with $\sigma_v \geq \sigma_{\min}$ in the current epoch may call \texttt{rotate(newRoot, manifestCid)}.
    \item \textbf{Slashing exposure}: a rotation whose root does not match the published manifest is provably false and triggers \texttt{slash(validatorId, amount)} against the rotator's stake.
    \item \textbf{Revenue share}: Treasury rewards for module-call fees are distributed proportional to staker STT, so honest publishers earn alongside the modules they index.
\end{enumerate}

\section{Priority Mechanism}
\label{sec:priority}

Within a single rotation the publisher batches many pending updates into one root. StakeTime priority settles two questions: which updates land in this epoch, and which modules occupy the 420-slot on-chain registry under contention.

\subsection{Per-update priority within an epoch}

Each pending update is signed by the module's owner and gossiped to a mempool. The publisher orders the mempool by:

\begin{equation}
\mathrm{priority}(u) = \mathrm{bond}(u) \cdot M(\ell_u) + \alpha \cdot \mathrm{age}(u)
\label{eq:priority}
\end{equation}

where $\mathrm{bond}(u)$ is the STT the submitter has locked behind the update, $M(\ell_u)$ is the same piecewise-linear lock multiplier StakeTime uses for stake weight (monotonic, on-chain enforced; see Equation~\ref{eq:multiplier-ref}), and $\alpha$ is a small fairness coefficient so old updates aren't permanently starved.

\begin{property}[Provable ordering]
The publisher MUST honour the priority order in Equation~\ref{eq:priority}. Any inversion is reconstructible from the mempool log---signed updates are deterministic---and triggers slashing.
\end{property}

\subsection{Slot priority for the competitive registry}

The full module set is unbounded off chain, but the \emph{on-chain discoverable} set is capped at 420 slots (StakeTime's Darwinian registry). When a 421st module would register, the slot held by the lowest-priority current occupant is forfeited:

\begin{equation}
\mathrm{slot\_score}(m) = \sum_{u \in \mathrm{users}} \mathrm{STT}_u \cdot w(m, u) \;-\; \pi \cdot \Delta_{\text{lastCall}}(m)
\label{eq:slot-score}
\end{equation}

where $w(m, u)$ is the usage weight (e.g.\ number of debit transactions a user has paid into $m$), $\pi$ is an inactivity-penalty rate, and $\Delta_{\text{lastCall}}$ is blocks since the last call.

\begin{itemize}
    \item New entries are \textbf{immune} from replacement for \texttt{immunityPeriod} blocks (default $\approx$ 1 day at 12 s).
    \item Registration costs a \textbf{locked bond} of governance tokens (default 1{,}000), forfeited on competitive replacement, returned on voluntary deregistration.
    \item The weakest occupant is found by \texttt{Registry.getWeakestSubnet()} in $O(1)$ via a heap maintained on each score update.
\end{itemize}

\begin{property}[Soft eviction]
Modules that fall out of the on-chain set remain fully accessible off chain via the manifest. They lose the gas-light on-chain discoverability hint until they earn it back through usage.
\end{property}

\subsection{Multiplier curve (reused from StakeTime Staking)}

\begin{equation}
M(x) = \begin{cases}
m_0 & x \leq b_0 \\[4pt]
m_i + \dfrac{(m_{i+1} - m_i)(x - b_i)}{b_{i+1} - b_i} & b_i < x \leq b_{i+1} \\[6pt]
m_{n-1} & x > b_{n-1}
\end{cases}
\label{eq:multiplier-ref}
\end{equation}

with strict monotonicity $b_i < b_{i+1}$ and weak monotonicity $m_i \leq m_{i+1}$ enforced on-chain to prevent arbitrage between lock durations.

\section{Gas and Storage Analysis}

\subsection{Per-operation gas}

\begin{table}[H]
\centering
\begin{tabularx}{\textwidth}{lXX}
\toprule
\textbf{Operation} & \textbf{Current Registry} & \textbf{Tree Registry} \\
\midrule
Register module          & 80{,}000--110{,}000 & $\approx 30{,}000 / N$ (amortised) \\
Update CID               & 30{,}000--50{,}000 & $\approx 30{,}000 / N$ (amortised) \\
Transfer ownership       & 25{,}000 & $\approx 30{,}000 / N$ \\
Revoke / delete          & 15{,}000 (SSTORE clear) & 25{,}000 (single SSTORE in \texttt{revoked} map) \\
Read by name             & free (view) & free (Merkle verify off-chain) \\
Enumerate all modules    & $O(N)$ event replay & 1 IPFS fetch of manifest \\
\bottomrule
\end{tabularx}
\caption{Gas comparison. $N$ is the rotation batch size (e.g.\ 100 updates per rotation).}
\end{table}

For $N = 100$, a rotation cycle costs $\approx 30{,}000 / 100 = 300$ gas per update---two orders of magnitude below the current floor.

\subsection{On-chain storage growth}

\begin{property}[Constant footprint]
The contract storage of \texttt{TreeRegistry} grows only with the number of actively revoked modules (a tiny minority). Compared to the linear $O(N)$ growth of the per-row registry, on-chain state is effectively constant.
\end{property}

\begin{property}[Sublinear emission cost]
A \texttt{Rotated} event is fixed-size regardless of batch contents. The \texttt{ModRegistered} event stream of the current design is $O(N)$ and quadratic to fully replay.
\end{property}

\section{Verifiability and Trust}

\subsection{What a client needs}

To verify that module \texttt{X} is registered to owner \texttt{O} with CID \texttt{C}:

\begin{enumerate}
    \item Fetch the current \texttt{root} from \texttt{TreeRegistry} (1 \texttt{eth\_call}).
    \item Fetch the manifest CID from the latest \texttt{Rotated} event (1 log filter).
    \item Pull the manifest from any IPFS gateway.
    \item Reconstruct the record locally and compute $\ell(r)$.
    \item Compute the Merkle proof against the on-chain root.
    \item (Optional) check that $\ell(r)$ is not in \texttt{revoked}.
\end{enumerate}

The client does \emph{not} need to trust the publisher: if the publisher signs a root that does not match the manifest, mismatched proofs reveal it instantly.

\begin{claim}[No fraudulent record can be proven]
\label{claim:no-fraud}
For any record $r' \notin \mathcal{T}$, no Merkle proof against $\mathrm{root}(\mathcal{T})$ verifies, under the standard collision-resistance assumption on \texttt{keccak256}.
\end{claim}

\subsection{Censorship}

The publisher can refuse to include a submitted update. We mitigate this two ways:

\begin{enumerate}
    \item \textbf{Forced inclusion.} Any user may pay gas to call \texttt{forceInclude(record, signature)} on a fallback contract that bypasses the publisher and triggers a special rotation. This restores the worst-case behaviour to today's per-row cost---but only when censored, not by default.
    \item \textbf{Publisher rotation.} If the publisher is a validator threshold (Phase 2), bribing $k$-of-$n$ validators is required to censor. Slashing is on the line.
\end{enumerate}

\section{Off-Chain Storage Architecture}

\subsection{Storage tiers}

\begin{table}[H]
\centering
\begin{tabularx}{\textwidth}{lX}
\toprule
\textbf{Tier} & \textbf{Contents} \\
\midrule
On-chain (1 slot) & Current root, epoch, publisher address \\
IPFS (manifests) & Full module records + tree structure, addressed by CID \\
Filecoin (long-term) & Manifest pinning with economic guarantees \\
Local cache & Compact LMDB of (leaf $\to$ proof) for each subscribed indexer \\
\bottomrule
\end{tabularx}
\caption{Storage tiering. The smaller the tier, the higher the trust.}
\end{table}

\subsection{Replication and availability}

A manifest is small (estimated $\approx 50$ bytes/record + tree overhead, so a 100k-module ecosystem is $\approx 5$ MB). Every running orbit node can pin every historical manifest at negligible cost. The protocol does not require a single trusted gateway; any IPFS or HTTP mirror suffices because authenticity is verified against the on-chain root.

\subsection{Module API surface}

This whitepaper ships as a runnable orbit module (\texttt{mod/orbit/whitepaper/}) with both a static document app and a working reference implementation of the tree:

\begin{lstlisting}[style=python]
m.fn('whitepaper/tree_build')(records)          # build tree + manifest
m.fn('whitepaper/tree_root')()                  # current root
m.fn('whitepaper/tree_proof')(name='agent')     # Merkle proof for a leaf
m.fn('whitepaper/tree_verify')(leaf, proof)     # client-side verify
\end{lstlisting}

\part{The Mod Protocol}

Part~I specifies the \emph{registry} layer---how module records are committed
and proven at scale. This part specifies what is being registered. The mod
protocol is a small set of conventions that turn any directory of code into a
\emph{module}: loadable from Python, callable from the CLI, servable as an HTTP
API, reachable by name through a gateway, legible to an agent, and---via
Part~I---anchored on chain. Everything else in the system is built from these
rules; at the time of writing the live fleet is 300+ orbit modules and a dozen
\texttt{core} modules, all obeying them.

\section{The Module}

\subsection{Anatomy}

The protocol's unit is the \emph{module}: a directory inside an \emph{orbit}
(\texttt{core} for shared infrastructure, \texttt{orbit} for applications,
plus \texttt{mods} for registry-installed and \texttt{local} for scratch),
containing an \emph{anchor file} and a \texttt{config.json}:

\begin{lstlisting}[style=python]
mod/orbit/<name>/
  mod.py         # anchor: class Mod - plain Python methods are the functions
  config.json    # name, description, owner, port, app_port, schema (CID)
  skill.md       # agent-facing docs (optional)
  README.md      # human docs (optional)
  src/ app/ ...  # anything else: Rust service, Next.js app, data
\end{lstlisting}

There is no decorator, manifest of exports, or registration step: the public
methods of the anchor class \emph{are} the module's functions, and the
framework derives their signatures by introspection
(\texttt{m.fns(mod)}, \texttt{m.schema(mod)}). An anchor named \texttt{mod.py},
\texttt{agent.py}, or \texttt{<name>.py} is equally valid; a module's
\emph{name} is derived from its path, so a directory is authoritative and two
orbits may not claim the same name.

\begin{lstlisting}[style=python]
{ "name": "whitepaper", "description": "...",
  "port": 50106, "app_port": 3106,
  "urls": { "api": "http://localhost:50106", "app": "http://localhost:3106" },
  "schema": "QmQQdHuUVasYCtHPrNBSNWmryzeDNVCpDavSu3rzpfuXEF" }
\end{lstlisting}

\texttt{config.json} is the module's public face and the exact object the
registry of Part~I hashes into a leaf. Its \texttt{schema} field is an IPFS CID
over the module's introspected function schema, which is what makes ``this
module still does what it claimed'' a checkable statement rather than a promise.

\subsection{Many implementations, one identity}

A module may ship any combination of implementations behind a single name: a
Python class, a compiled service API (Rust/axum is common), and a web app
(Next.js). The protocol cares about a module's \emph{name}, \emph{owner}, and
\emph{endpoints}, not its language. This module is itself an example: the same
Merkle tree is implemented in Python (\texttt{mod.py}) and in Rust
(\texttt{src/api}), both reading and writing the same state file, both applying
the same canonical-JSON rules so that the roots they produce are identical.

Cross-cutting capabilities---identity, process supervision, environment,
storage, routing, registry---are themselves modules under \texttt{core/} and are
reused rather than re-implemented. This is what makes the system modular all the
way down: the protocol is written in its own units.

\section{Calling a Module}

\subsection{Three surfaces, one semantics}

A function is invoked identically from three places:

\begin{lstlisting}[style=python]
m.fn('whitepaper/tree_proof')(name='agent')       # Python
m whitepaper/tree_proof name=agent                # CLI
POST /api/whitepaper/tree/proof  {"name":"agent"} # HTTP
\end{lstlisting}

Serving is not a rewrite: \texttt{m serve <name>} places the module behind the
core server and every public function becomes a POST endpoint, with JSON
arguments in the body and \texttt{\{"result": ...\}} in the response. Modules
with native servers (Rust, Node, Next.js) expose the same public shape---an API
on \texttt{port}, an app on \texttt{app\_port}---and are therefore
indistinguishable to a caller.

\subsection{The null call}

A POST to a module's bare URL, with no function named, returns the module's
\texttt{info}: its name, its functions, and its schema.

\begin{lstlisting}[style=python]
POST /mod/<name>      ->  { name, functions, schema, ... }
\end{lstlisting}

This single rule makes the ecosystem self-describing without a service
directory: any endpoint that speaks the protocol can be asked what it can do,
and the answer is derived from the code rather than maintained beside it. It is
also the bridge to Part~I---the schema a module reports is the schema whose CID
its registry leaf commits to.

\section{Addressing: One Name, Three Forms}

Every deployment routes modules by the same rule. A module named \texttt{X} is
reachable as:

\begin{table}[H]
\centering
\begin{tabularx}{\textwidth}{llX}
\toprule
\textbf{Form} & \textbf{Target} & \textbf{Rule} \\
\midrule
\texttt{https://<host>/X} & app (\texttt{app\_port}) & prefix \emph{kept}; the app sets \texttt{basePath /X} \\
\texttt{https://<host>/api/X} & API (\texttt{port}) & prefix \emph{stripped}: \texttt{/api/X/y} $\to$ \texttt{/y} \\
\texttt{https://X.<host>} & app & name form, resolved by the DNS layer \\
\bottomrule
\end{tabularx}
\caption{The URL rule. App at \texttt{/X}, API at \texttt{/api/X}, discovery by null call.}
\end{table}

Three interchangeable implementations enforce it: the production Caddy gateway,
whose routes are \emph{generated} from each module's \texttt{config.json} rather
than hand-written; \texttt{routy}, a standalone Rust gateway reserving
\texttt{/\_api/*} for its own control plane so that no module name is shadowed;
and a Next.js middleware applying the identical rewrite inside the core
frontend. Because the rule is mechanical, a module's routing is a consequence of
its config, not a deployment artifact---which is precisely the property that
lets a registry leaf, rather than an operator, determine where a name resolves.

The \texttt{dns} module closes the loop for the third form: it derives an
authoritative zone from the module fleet and serves it over UDP and TCP, so
\texttt{X.<host>}, \texttt{<host>/X}, and \texttt{<host>/api/X} agree by
construction rather than by convention. Attribution records (\texttt{\_mod} TXT)
carry the declared owner and CID of the module answering at a name, making the
DNS layer a second, independent witness to the same facts the tree commits.

\section{Runtime and Environment}

\subsection{One environment, defined once}

Historically each module shipped its own \texttt{Dockerfile}, re-deriving Node,
Python, Rust, and a reverse proxy per image. The protocol replaces this with a
single reusable environment, \texttt{core/nix} (``modenv''), pinned via a Nix
flake. \texttt{packages.nix} is the single source of truth---language toolchains
and the system tools the fleet needs, but \emph{not} application libraries
(those stay with \texttt{pip}/\texttt{npm}/\texttt{cargo}), keeping the shared
layer small. A module opts in by shipping a two-line flake that reuses
\texttt{modenv.lib.mkShell}; otherwise it inherits the shared environment. This
is to environments what a shared supervisor is to processes: define once, use
everywhere.

\subsection{Supervision: core/pm}

Process supervision is unified in \texttt{core/pm}, which launches every service
inside its \emph{nix image}---the module's own flake if present, else the shared
\texttt{core/nix}. The mechanism imports the image's environment and then
\texttt{exec}s the real server, so the supervisor tracks the actual process:

\begin{lstlisting}[style=python]
cd <service dir>
eval "$(nix print-dev-env <image>)"   # import the pinned toolchain
exec bash start.sh                     # become the server; supervisor tracks it
\end{lstlisting}

The supervisor backend (today \texttt{pm2}) is pluggable. The result is
hermetic, reproducible execution without per-module containers: the apt-based
image becomes a pinned Nix image, and the supervisor is shared.

\section{Gateway and On-Demand Execution}

\subsection{Federated front door}

All traffic enters through a single gateway. Some modules are routed directly by
host rules; the remainder are federated through a portal that maps
\texttt{/\{mod\}} and \texttt{/api/\{mod\}} to each module's registered
\texttt{app}/\texttt{api} endpoints, resolved from the registry's app namespace.
A request thus reaches a module by \emph{name}, independent of where it runs---
the same indirection that lets the registry of Part~I move a name to a new host
or a new owner without any client learning a new address.

\subsection{Scale-to-zero}

An \emph{activator} sits in front of the local fleet. It stops any module idle
beyond a threshold with no active connections, and \emph{wakes} a stopped module
on the next request---starting its supervised processes, waiting for the port,
then proxying---so a slept module costs nothing yet remains reachable. Because
the host owns the machine, manual control is authoritative over the automation:
modules may be \emph{disabled} (kept down; the activator refuses to wake them)
or \emph{pinned} (never slept), via host-local overrides. This gives the
ecosystem the elasticity of serverless execution over a fleet of long-lived
modules, and it is why a 300-module fleet fits on one machine: the marginal cost
of a registered-but-idle module is zero.

\section{Identity and Authorization}

\subsection{Keys}

Identity is a keypair, not an account. \texttt{core/key} manages multi-chain
keys---ECDSA (Ethereum), sr25519 and ed25519 (Substrate), Solana---stored
encrypted under \texttt{\textasciitilde/.mod/key/}. The same key material that
signs a registry update in Part~I signs an HTTP request in Part~II; there is no
second identity system to reconcile.

\subsection{The token}

\texttt{core/server/auth} (\texttt{m.mod('auth')}) is the one identity layer
modules share. A token is a signature over a timestamped envelope, base64url
encoded:

\begin{lstlisting}[style=python]
token = b64url({ "data": <payload>,
                 "time": "<unix seconds>",
                 "key":  "<signer address>",
                 "signature": sign(hash(data|time|key)) })
\end{lstlisting}

Verification recovers the signer from the signature, checks it against the
declared \texttt{key}, and rejects the token if \texttt{time} is older than the
configured maximum age---so a captured token is not a durable credential. The
crypto type is inferred from the address format, and browser wallets are
first-class: a \texttt{personal\_sign} over the same envelope mints a valid mod
token, with legacy \texttt{v=27/28} recovery identifiers normalised on verify.
No module rolls its own scheme; ``who is calling'' is answered the same way
everywhere.

\subsection{The owner/peer model}

Authorization is a two-tier model layered on that identity:

\begin{itemize}
    \item \textbf{Owner.} The address in a module's \texttt{config.json}
    (optionally extended to on-chain holders). The owner has authority over the
    module---and, for the host operator, over the module tree. Owner-only
    powers (process control, deletion, managing the sign-in list) never widen
    beyond this.
    \item \textbf{Peer.} Every other authenticated identity. A peer is confined
    to its own off-chain workspace,
    \texttt{\textasciitilde/.mod/peers/<addr>/}, and its jobs are
    privilege-dropped---kernel-enforced containment that no prompt or path can
    escape. A peer can build and run \emph{their} modules without ever reaching
    the owner's tree.
\end{itemize}

An optional off-chain whitelist gates \emph{sign-in} but confers no ownership; a
whitelisted address is still a peer for editing. All such private authorization
state lives off-chain (under \texttt{\textasciitilde/.mod/}), never in the
committed, publicly mirrored \texttt{config.json}---so the public tree of Part~I
carries only public facts, while access control stays private to the host. This
separation is load-bearing for the registry: a world-readable manifest is safe
precisely because nothing secret is ever a leaf.

\section{State, Storage, and Content Addressing}

Module state persists as JSON under \texttt{\textasciitilde/.mod/<module>/},
optionally AES-encrypted, and is deliberately outside the module directory: the
code is public and mirrorable, the state is not. \texttt{core/store} adds a
content-addressed layer on top---objects are written and retrieved by CID, so a
shared artifact (a schema, a manifest, a build) is named by its hash rather than
by its location.

This is the same addressing discipline the registry uses, applied one level
down, and it is what makes the three layers compose:

\begin{table}[H]
\centering
\begin{tabularx}{\textwidth}{lX}
\toprule
\textbf{Layer} & \textbf{Names things by} \\
\midrule
Chain (Part~I) & one root---a commitment to the whole set \\
Content store & CID---a commitment to one object \\
Filesystem / DNS & name---a mutable pointer, checkable against both above \\
\bottomrule
\end{tabularx}
\caption{Three naming layers. Mutability increases downward; each layer is
verifiable against the one above it.}
\end{table}

\section{The Agent Interface}

The protocol was designed on the assumption that most callers would eventually
be programs, not people. Three properties make a module legible to an agent
without any per-module integration work:

\begin{enumerate}
    \item \textbf{Introspection.} Function names and signatures are derived from
    the code, so a tool list can be generated rather than maintained.
    \item \textbf{\texttt{skill.md}.} An agent-facing document beside the
    human-facing \texttt{README.md}, stating what the module is for and how to
    drive it. Two audiences, two files, one module.
    \item \textbf{MCP.} A module exposes its functions as Model Context Protocol
    tools from one method---the same functions the CLI and HTTP surfaces
    expose---so an agent connects to a module the same way a browser does.
\end{enumerate}

Agentic publishing is therefore not a special case bolted onto the registry: an
agent that writes a directory with an anchor class and a \texttt{config.json}
has produced a module that is immediately callable, routable, and registrable by
the same rules as a hand-written one. This is the growth curve Part~I is sized
for, and the reason per-row registration cost is the binding constraint rather
than a nuisance.

\section{Protocol Invariants}

The whole protocol reduces to a handful of rules. Every module in the fleet
satisfies all of them; a directory that satisfies all of them is a module.

\begin{table}[H]
\centering
\begin{tabularx}{\textwidth}{lX}
\toprule
\textbf{Invariant} & \textbf{Statement} \\
\midrule
Definition & A module is a directory with an anchor class and a \texttt{config.json}. \\
Exposure & The public methods of the anchor class are the module's functions. \\
Equivalence & Python, CLI, and HTTP call the same function with the same semantics. \\
Discovery & A null call to a module returns its name, functions, and schema. \\
Addressing & App at \texttt{/\{mod\}}, API at \texttt{/api/\{mod\}}, name at \texttt{\{mod\}.\{host\}}. \\
Identity & One shared auth module; a token is a time-bounded signature. \\
Authority & Owner from \texttt{config.json}; everyone else is a confined peer. \\
Privacy & Public code and config are committed; state and ACLs stay under \texttt{\textasciitilde/.mod/}. \\
Provenance & \texttt{config.json} is the registry leaf; its \texttt{schema} CID commits to the function surface. \\
\bottomrule
\end{tabularx}
\caption{The mod protocol in nine rules.}
\end{table}

\part{Adoption}

\section{Migration Path}

\begin{enumerate}
    \item \textbf{Shadow mode.} \texttt{TreeRegistry} deploys alongside the existing per-row \texttt{Registry}. The publisher mirrors all \texttt{Registry} events into the tree. Both registries serve simultaneously.
    \item \textbf{Indexer flip.} Off-chain indexers (the \texttt{api} module, the polymarket frontend, ecosystem block explorers) cut their reads over to the tree-based view. Writes still go to the legacy registry.
    \item \textbf{Write deprecation.} New \texttt{registerMod} calls are routed through a meta-aggregator that submits to the tree pool. The legacy contract is frozen but its state is mirrored permanently into the tree's genesis manifest.
    \item \textbf{Sunset.} The legacy registry stays readable; no new writes are accepted. Eventually it is marked \texttt{setOwnerless} and forgotten.
\end{enumerate}

\section{Comparison to Related Work}

\begin{itemize}
    \item \textbf{Ethereum Name Service (Namewrapper).} ENS stores per-name records on-chain. Our design is closer to \emph{rollup} commit roots applied to a developer registry.
    \item \textbf{Optimistic rollups.} We borrow the publish-root-then-prove-fraud pattern, but without fraud proofs: there is no state transition to dispute, only a static set membership. This makes the model significantly simpler.
    \item \textbf{IPNS / IPFS DNSLink.} Both lack an on-chain anchor, so they cannot resolve disputes about which manifest is canonical. Our root + epoch counter solves that.
    \item \textbf{Sparse Merkle Trees (SMT).} An SMT would give us inclusion + non-inclusion proofs for free at the cost of larger proofs. For an opt-in publisher model where revocation is a tombstone bit, we do not need non-inclusion proofs; the sorted-leaf binary tree is sufficient and cheaper.
\end{itemize}

\section{Conclusion}

Storing the tree instead of the leaves shifts the cost curve from $O(N)$ to $O(1)$ on chain, eliminates per-module gas as a publishing tax, and makes the registry trivially mirrorable. It preserves authenticity through Merkle proofs and ownership through signed updates, and degrades gracefully under censorship via forced inclusion. The resulting system is the registry layer the MOD orbit ecosystem needs to grow from hundreds of modules to millions, including those published autonomously by agents.

That layer is only useful because of what it registers. The mod protocol of Part~II makes a module a directory, its anchor class its API, its \texttt{config.json} its public face, and its name a route---so a registry leaf is not a description of a module but the module's own self-description, mechanically derived and mechanically checkable. The chain commits to a root, the content store commits to each object, and the filesystem and DNS hold the mutable pointers that both can audit. A module written by an agent this afternoon obeys the same nine invariants as \texttt{core/auth}, is reachable at the same three URLs, and enters the same tree. Scaling the registry and specifying the protocol are therefore one problem: the cheaper a name is to commit, the more of the ecosystem can afford to be verifiable.

\appendix

\section{Reference Implementation}

The module ships a Python reference implementation at \texttt{mod/orbit/whitepaper/mod.py} exposing:

\begin{lstlisting}[style=python]
class Mod:
    def tree_build(self, records: list[dict]) -> dict: ...
    def tree_root(self) -> str: ...
    def tree_proof(self, name: str) -> list[str]: ...
    def tree_verify(self, leaf: str, proof: list[str]) -> bool: ...
\end{lstlisting}

\section{Notation}

\begin{tabularx}{\textwidth}{lX}
\toprule
$N$ & Total number of modules in the registry. \\
$r$ & A module record (canonical JSON). \\
$\ell(r)$ & The leaf hash $\mathrm{keccak256}(\mathrm{cjson}(r))$. \\
$\mathcal{T}$ & The Merkle tree over all leaves. \\
$R$ & The root of $\mathcal{T}$, anchored on-chain. \\
$E$ & The epoch counter. \\
\bottomrule
\end{tabularx}

\end{document}
